\documentclass[10pt, conference, compsocconf]{IEEEtran}

\usepackage{color}
\usepackage{soul}

%\documentclass[a4paper, UKenglish]{lipics}
%\usepackage{amsmath}
%\theoremstyle{plain}

\usepackage[noend]{algorithm2e}
\usepackage{algorithmicx}
\usepackage{graphicx}
\usepackage{amssymb}
\usepackage{amsmath}
\usepackage{times}
\let\proof\relax % undefine the environment
\let\endproof\relax
\usepackage{amsthm}

\newcommand{\todo}[1]{\color{red}\textbf{\hl{\{ #1 \}}}\color{black}\xspace}


\begin{document}
\title{High Performance Space-Time Kernel Density Estimation}


\author{\IEEEauthorblockN{Dinesh Panchananam$^\dag$, Alex Hohl$^\ddag$, Wenwu Tang$^\ddag$, Eric Delmelle$^\ddag$, Erik  Saule$^\dag$}
\IEEEauthorblockA{$^\dag$ Dept. of Computer Science, $^\ddag$Dept. of Geography and Earth Sciences\\
UNC Charlotte\\
Charlotte, NC, USA\\
Email: {dpanchan,ahohl,wtang4,eric.delmelle,esaule}@uncc.edu}
}

\maketitle


\begin{abstract}
  The exponential growth of available data, caused by the BigData era,
  increased the interest in interactive exploratory analysis as dataset
  can no longer be understood by crawling manually through it and
  looking at a few statistics. In a Geographical Information Systems
  (GIS), the dataset is often composed of events localized in space
  and time; and visualizing such a dataset involves building a map of
  where the events occurred.

  We focus in this paper on events that are localized among three
  dimensions (latitude, longitude, and time), and on computing the
  first step of the visualization pipeline, space-time kernel density
  estimation, which is most computationally expensive.  Starting from
  a gold standard implementation, we show how algorithm design and
  engineering, parallel decomposition, and acceleration can be applied
  to bring real-time computing to space-time kernel density
  estimation. We validate our techniques on real world datasets
  extracted from infectious disease, social media, and ornithology.
\end{abstract}

\begin{IEEEkeywords}
space-time kernel density; performance; distributed memory; gpu
\end{IEEEkeywords}


\IEEEpeerreviewmaketitle

\section{Introduction}

\section{Space-Time Kernel Density Estimation}

\subsection{Problem Definition}

\subsection{Gold Standard Implementation}

\subsection{Related Works}

\section{Algorithm Design and Engineering}

\subsection{Observation-based Algorithm}

\subsection{Exploiting Symmetries}

\section{Utilizing Modern Hardware}
\todo{That's a terrible section header...}

\subsection{Spatial Decomposition}

\subsection{Observation Decomposition}

\subsection{Vectorization}

\subsection{Accelerators}


\section{Experiments}

\subsection{Experimental Setting}

\paragraph{Machine}

\paragraph{Dataset}

\subsection{Results}

\begin{table*}
  \begin{tabular}{|l|rrrr|r|}
Instance & PB & PB-SYMDISK & PB-SYMBAR & PB-SYM & SYM speedup \\
Dengue\_lowres-lowbw    & 0.031         & 0.027         & 0.028         & 0.027         & 1.148 \\
Dengue\_lowres-highbw   & 0.856         & 0.504         & 0.755         & 0.503         & 1.702 \\
Dengue\_highres-lowbw   & 0.082         & 0.080         & 0.080         & 0.080         & 1.025 \\
Dengue\_highres-highbw  & 3.437         & 2.044         & 3.026         & 2.087         & 1.647 \\
Dengue\_highres-vhighbw         & 49.585        & 12.002        & 40.791        & 7.424         & 6.679 \\
Pollen\_lowres-lowbw    & 0.619         & 0.226         & 0.498         & 0.210         & 2.948 \\
Pollen\_highres-lowbw   & 19.183        & 7.268         & 15.133        & 5.210         & 3.682 \\
Pollen\_highres-medbw   & 330.352       & 87.250        & 270.699       & 56.274        & 5.870 \\
Pollen\_highres-highbw  & 2596.372      & 626.069       & 2135.457      & 394.337       & 6.584 \\
Flu-Animal\_lowres-lowbw        & 0.030         & 0.030         & 0.030         & 0.031         & 0.968 \\
Flu-Animal\_lowres-highbw       & 0.058         & 0.040         & 0.051         & 0.039         & 1.487 \\
Flu-Animal\_medres-lowbw        & 0.287         & 0.272         & 0.283         & 0.273         & 1.051 \\
Flu-Animal\_medres-highbw       & 0.604         & 0.364         & 0.527         & 0.331         & 1.825 \\
Flu-Animal\_highres-lowbw       & 7.142         & 6.737         & 7.019         & 5.542         & 1.289 \\
Flu-Animal\_highres-highbw      & 12.086        & 6.817         & 10.448        & 6.401         & 1.888 \\
  \end{tabular}
  \end{table*}

\section{Conclusion}

\section*{Acknowledgment}





\bibliographystyle{IEEEtran}
\bibliography{ipdps}




\end{document}


